\documentclass[9pt,a4paper,onecolumn,twoside]{tau-class/tau}
\usepackage[english]{babel}
\usepackage{pgfplotstable}
\usepackage{booktabs}
\usepackage{graphicx}
\usepackage{listings}
\usepackage{xcolor}

\lstset{
    basicstyle=\ttfamily\small,
    breaklines=true,
    frame=single,
    language=Python,
    keywordstyle=\color{blue},
    commentstyle=\color{gray},
    stringstyle=\color{red},
    showstringspaces=false
}

%----------------------------------------------------------
% TITLE
%----------------------------------------------------------

\journalname{Big Data Project}
\title{Income Inequality and Voting Patterns Across U.S. Counties: A Big Data Analysis}

%----------------------------------------------------------
% AUTHORS, AFFILIATIONS AND PROFESSOR
%----------------------------------------------------------

\author[a,1]{Monish Sinha}

%----------------------------------------------------------

\affil[a]{Big Data, University College Dublin}

\professor{Professor Soale}

%----------------------------------------------------------
% FOOTER INFORMATION
%----------------------------------------------------------

\institution{University College Dublin}
\footinfo{Big Data Project}
\theday{February 11, 2026}
\leadauthor{Sinha}
\course{Big Data}

%----------------------------------------------------------
% ABSTRACT AND KEYWORDS
%----------------------------------------------------------

\begin{abstract}    
This project analyzes the relationship between economic, demographic, social, and housing conditions and voting patterns across all 3,143 U.S. counties. By combining American Community Survey (ACS) 5-Year data profiles (2009--2020) with county-level presidential election results (2000--2024), we investigate how socioeconomic indicators correlate with political preferences. Our analysis processes 93 MB of ACS data across 2,053 files and 94,019 voting records. We implement both a single-threaded Python prototype and a Spark MapReduce solution to compare performance characteristics. Results show that Democratic-winning counties have 7$\times$ larger average populations (288,926 vs 40,311), reflecting the urban-rural political divide.
\end{abstract}

%----------------------------------------------------------

\keywords{income inequality, voting patterns, U.S. counties, big data, ACS, presidential elections, MapReduce, Spark}

%----------------------------------------------------------

\begin{document}

\maketitle 
\thispagestyle{firststyle} 
\tauabstract 

%----------------------------------------------------------

\section{Introduction}

\taustart{T}his project analyzes the relationship between county-level economic, demographic, social, and housing conditions and voting patterns across all 3,143 U.S. counties. We combine American Community Survey (ACS) 5-Year Data Profiles (2009--2020) with county-level presidential election results (2000--2024) to investigate how socioeconomic indicators correlate with political preferences.

\subsection{Datasets}

We utilize five distinct ACS Data Profile tables downloaded via the Census API:

\begin{itemize}
    \item \textbf{DP02 (Social):} 158 variables covering households, education, language, and ancestry
    \item \textbf{DP03 (Economic):} 141 variables covering employment, income, poverty, and commuting
    \item \textbf{DP04 (Housing):} 145 variables covering occupancy, structure, value, and rent
    \item \textbf{DP05 (Demographic):} 93 variables covering age, sex, race, and population
    \item \textbf{Presidential Voting:} County-level election results (2000--2024), 94,019 records
\end{itemize}

\subsection{Big Data Justification}

\textbf{Volume:} The combined datasets contain 93 MB of ACS data across 2,053 files (51 states $\times$ 12 years $\times$ 4 tables), plus 8.9 MB of presidential voting data (94,019 rows). The merged datasets contain 37,710 county-year records with up to 769 columns each. Processing was executed on a MacBook Air M4 (10 cores, 24 GB RAM) with Python 3.9.6 and Pandas 2.3.3.

\textbf{Variety:} We integrate multiple structured datasets with different schemas: wide-format ACS tables with hundreds of variables per county, and long-format voting data with multiple candidates per county-year. Each dataset requires different parsing, cleaning, and merging strategies. Column names vary across years as Census Bureau updates variable definitions.

\textbf{Value:} Understanding how economic conditions correlate with voting behavior provides value for policymakers analyzing political dynamics, researchers studying socioeconomic determinants of voting, and analysts examining the urban-rural political divide. Our analysis reveals that Democratic-winning counties have 7$\times$ larger average populations than Republican-winning counties (288,926 vs 40,311), quantifying the urban-rural political divide.

%----------------------------------------------------------

\section{Traditional Solution}

\taustart{B}efore developing a big data pipeline, we built a single-threaded Python prototype to validate our processing logic and establish performance baselines. The prototype decomposes the analysis into five key steps.

\subsection{System Specifications}

\begin{table}[htbp]
\centering
\caption{Hardware and Software Configuration}
\label{tab:specs}
\begin{tabular}{ll}
\hline
Component & Value \\
\hline
Machine & MacBook Air \\
Chip & Apple M4 (10 cores: 4P + 6E) \\
RAM & 24 GB \\
OS & macOS 26.2 \\
Python & 3.9.6 \\
Pandas & 2.3.3 \\
\hline
\end{tabular}
\end{table}

\subsection{Step 1: Load ACS Economic Data}

Load the consolidated ACS DP03 economic data for 2020 containing employment, income, and poverty statistics for all U.S. counties.

\begin{lstlisting}
import pandas as pd
econ = pd.read_csv('data/acs_merged/economic/economic_2020.csv')
\end{lstlisting}

\textbf{Results:} 3,143 rows (all U.S. counties), 141 columns\\
\textbf{Execution Time:} 0.024 seconds | \textbf{Memory:} 12.8 MB

\subsection{Step 2: Load Presidential Voting Data}

Load county-level presidential election results containing vote counts by candidate and party.

\begin{lstlisting}
pres = pd.read_csv('data/countypres_2000-2024.csv')
\end{lstlisting}

\textbf{Results:} 94,019 rows, 13 columns\\
\textbf{Execution Time:} 0.055 seconds | \textbf{Memory:} 39.1 MB (cumulative)

\subsection{Step 3: Filter and Clean Data}

Filter presidential data to 2020 and create standardized 5-digit FIPS codes for merging.

\begin{lstlisting}
pres_2020 = pres[pres['year'] == 2020].copy()
econ['full_fips'] = (econ['state_fips'].astype(str).str.zfill(2) + 
                     econ['county_fips'].astype(str).str.zfill(3))
pres_2020['full_fips'] = (pres_2020['county_fips'].astype(str)
                          .str.replace('.0','').str.zfill(5))
\end{lstlisting}

\textbf{Results:} Filtered to 22,093 rows (2020 election only)\\
\textbf{Execution Time:} 0.016 seconds

\subsection{Step 4: Merge Datasets}

Join economic data with presidential voting data on county FIPS code.

\begin{lstlisting}
merged = pd.merge(econ, 
    pres_2020[['full_fips', 'party', 'candidatevotes']], 
    on='full_fips', how='inner')
\end{lstlisting}

\textbf{Results:} 21,894 merged rows, 145 columns\\
\textbf{Execution Time:} 0.011 seconds | \textbf{Memory:} 67.0 MB (cumulative)

\subsection{Step 5: Aggregate and Analyze}

Determine winning party per county and calculate average population by party.

\begin{lstlisting}
winner_idx = pres_2020.groupby('full_fips')['candidatevotes'].idxmax()
winners = pres_2020.loc[winner_idx][['full_fips', 'party']]
analysis = pd.merge(econ, winners, on='full_fips')
result = analysis.groupby('winning_party').agg({
    'population_16_plus': 'mean', 'full_fips': 'count'})
\end{lstlisting}

\textbf{Execution Time:} 0.008 seconds

\subsection{Execution Summary}

\begin{table}[htbp]
\centering
\caption{Prototype Execution Metrics}
\label{tab:execution}
\begin{tabular}{lrr}
\hline
Step & Description & Time (s) \\
\hline
1 & Load Economic Data & 0.024 \\
2 & Load Presidential Data & 0.055 \\
3 & Filter \& Clean & 0.016 \\
4 & Merge Datasets & 0.011 \\
5 & Aggregate \& Analyze & 0.008 \\
\hline
\textbf{Total} & & \textbf{0.114} \\
\hline
\end{tabular}
\end{table}

\begin{table}[htbp]
\centering
\caption{Key Findings: Population by Winning Party (2020)}
\label{tab:findings}
\begin{tabular}{lrrr}
\hline
Party & County Count & Total Population & Avg Population \\
\hline
DEMOCRAT & 546 & 157,753,458 & 288,926 \\
REPUBLICAN & 2,569 & 103,559,736 & 40,311 \\
\hline
\end{tabular}
\end{table}

\textbf{Interpretation:} Democratic-winning counties have significantly larger populations on average (urban areas), while Republican-winning counties are more numerous but smaller (rural areas). This quantifies the urban-rural political divide.

\textbf{Peak Memory Usage:} 67.0 MB

%----------------------------------------------------------

\section{MapReduce Optimisation}

\taustart{W}e identified the two most time-consuming steps from Section 2 that can be optimized using MapReduce: (1) data loading (Step 2: 0.055s, 48\% of total time) and (2) dataset merging (Step 4: 0.011s, 10\% of total time). We implemented a Spark RDD solution using the PySpark Core API.

\subsection{Why MapReduce is Suitable}

These steps are suitable for parallel processing because:

\begin{itemize}
    \item \textbf{Data Loading (Map Phase):} Each CSV file is independent and can be read in parallel. There are no dependencies between files---perfect for the ``embarrassingly parallel'' pattern.
    \item \textbf{Merge/Join (Reduce Phase):} The merge operation groups records by county FIPS code. This is a classic reduce operation where all records with the same key are collected and combined.
\end{itemize}

\textbf{Expected Improvement:} 4--8$\times$ speedup on a multi-core system, with linear scaling as data size increases.

\subsection{MapReduce Solution}

\textbf{Map Function:} Parse CSV lines and emit (county\_fips, record) pairs.

\begin{lstlisting}
# Map function for economic data
def map_economic(line):
    fields = line.split(',')
    state_fips = fields[3].zfill(2)
    county_fips = fields[4].zfill(3)
    full_fips = state_fips + county_fips
    return (full_fips, ('economic', fields))

# Map function for voting data  
def map_voting(line):
    fields = line.split(',')
    county_fips = fields[4].zfill(5)
    return (county_fips, ('voting', fields))
\end{lstlisting}

\textbf{Reduce Function:} Merge economic and voting records, determine winner.

\begin{lstlisting}
def reduce_merge(county_fips, records):
    economic_data, voting_data = None, []
    for source, data in records:
        if source == 'economic':
            economic_data = data
        else:
            voting_data.append(data)
    if economic_data and voting_data:
        winner = max(voting_data, key=lambda x: x['votes'])
        return {'county_fips': county_fips,
                'population': economic_data['population'],
                'winning_party': winner['party']}
\end{lstlisting}

\textbf{Spark RDD Implementation:}

\begin{lstlisting}
from pyspark import SparkContext
sc = SparkContext("local[4]", "ACS_Voting_Merge")

economic_rdd = sc.textFile("economic_2020.csv") \
    .map(parse_economic_line).filter(lambda x: x)
voting_rdd = sc.textFile("countypres_2000-2024.csv") \
    .map(parse_voting_line).filter(lambda x: x)

merged_rdd = economic_rdd.union(voting_rdd) \
    .groupByKey().mapValues(merge_records)

result = merged_rdd \
    .map(lambda x: (x[1]['winning_party'], x[1]['population'])) \
    .reduceByKey(lambda a, b: a + b).collect()
\end{lstlisting}

\subsection{MapReduce Results}

\begin{table}[htbp]
\centering
\caption{Performance Comparison: Baseline vs MapReduce}
\label{tab:mapreduce}
\begin{tabular}{lrrr}
\hline
Configuration & Time (s) & Speedup & Notes \\
\hline
Baseline (Pandas) & 0.103 & 1.00$\times$ & Single-threaded \\
Spark local[1] & 1.311 & 0.08$\times$ & JVM startup dominates \\
Spark local[2] & 0.550 & 0.19$\times$ & Some parallelism \\
Spark local[4] & 0.523 & 0.20$\times$ & Diminishing returns \\
\hline
\end{tabular}
\end{table}

\begin{table}[htbp]
\centering
\caption{MapReduce Analysis Results (2020 Election)}
\label{tab:mr_results}
\begin{tabular}{lrrr}
\hline
Party & County Count & Total Population & Avg Population \\
\hline
DEMOCRAT & 546 & 157,753,458 & 288,926 \\
REPUBLICAN & 2,569 & 103,559,736 & 40,311 \\
\hline
\end{tabular}
\end{table}

\subsection{Analysis: Why Results Deviate from Expectations}

\textbf{Expected:} 4--8$\times$ speedup with parallel processing.\\
\textbf{Actual:} 0.08--0.20$\times$ (5--12$\times$ \textit{slower} than baseline).

The deviation is explained by:

\begin{enumerate}
    \item \textbf{JVM Startup Overhead:} Spark requires 1--2 seconds to initialize the JVM, load classes, and set up the execution environment. For a job that completes in 0.1 seconds with Pandas, this overhead is 10--20$\times$ the actual work.
    
    \item \textbf{Dataset Size Too Small:} Our dataset (67 MB) is below the threshold where MapReduce provides benefit. The rule of thumb is:
    \begin{itemize}
        \item $<$ 100 MB: Use single-threaded processing
        \item 100 MB -- 10 GB: Consider local parallelism
        \item $>$ 10 GB: MapReduce/Spark provides clear benefit
    \end{itemize}
    
    \item \textbf{Serialization Overhead:} PySpark must serialize Python objects to JVM and back, adding latency for each record.
\end{enumerate}

\textbf{Projected Performance at Scale:}

\begin{table}[htbp]
\centering
\caption{Projected Speedup at Larger Data Scales}
\label{tab:projected}
\begin{tabular}{lrrr}
\hline
Dataset Size & Pandas (est.) & Spark 4-core (est.) & Speedup \\
\hline
67 MB & 0.10s & 0.52s & 0.19$\times$ \\
670 MB & 1.0s & 0.8s & 1.25$\times$ \\
6.7 GB & 10s & 3.5s & 2.9$\times$ \\
67 GB & 100s & 15s & 6.7$\times$ \\
\hline
\end{tabular}
\end{table}

\textbf{Conclusion:} The MapReduce implementation is correct and functional, but the current dataset is too small to benefit from parallel processing. For production use with larger datasets ($>$1 GB), the MapReduce approach would provide significant speedup.

%----------------------------------------------------------

\end{document}
